\subsection{Цель \texttt{report}}

\subsubsection{Листинг}

\begin{lstlisting}[caption={Часть содержимого файла \texttt{build.xml}},captionpos=b]
    <target name="report">
        <antcall target="test">
          <param name="fse.test.report.output" value="true" />
        </antcall>

        <echo message="Saving test reports to SVN..."/>

        <java classname="org.tmatesoft.svn.cli.SVN"
              classpathref="fse.ivy.path.ant"
              fork="true">
          <arg value="update"/>
        </java>

        <java classname="org.tmatesoft.svn.cli.SVN"
              classpathref="fse.ivy.path.ant"
              fork="true">
            <arg value="add"/>
            <arg value="--depth"/>
            <arg value="empty"/>
            <arg value="--force"/>
            <arg value="${fse.test.target}"/>
        </java>

        <java classname="org.tmatesoft.svn.cli.SVN"
              classpathref="fse.ivy.path.ant"
              fork="true">
            <arg value="add"/>
            <arg value="--force"/>
            <arg value="${fse.test.target}/${fse.test.reports.subdir}"/>
        </java>

        <java classname="org.tmatesoft.svn.cli.SVN"
              failonerror="true"
              classpathref="fse.ivy.path.ant"
              fork="true">
            <arg value="commit"/>
            <arg value="-m"/>
            <arg value="feat: add test reports (Ant-automated)"/>
        </java>

        <echo message="Done task 'report'!" />
    </target>
\end{lstlisting}

\subsubsection{Комментарии}

Сохраняет результаты тестирования в SVN-репозиторий. Сначала вызывает цель \texttt{test}, чтобы гарантировать генерацию свежих XML-отчётов. Затем выполняет SVN-операции: обновление рабочей копии, добавление директории \texttt{build-test} (только структуры, без содержимого) и добавление поддиректории с отчётами. После этого выполняет коммит с сообщением. Все операции используют \texttt{SVNKit} через задачу java. Если SVN недоступен или директория не является рабочей копией, цель завершается с ошибкой.

\subsubsection{Значения используемых параметров}

\begin{lstlisting}[language=Properties,caption={Часть содержимого файла \texttt{build.properties}},captionpos=b]
fse.test.target = build-test
fse.test.reports.subdir = reports
\end{lstlisting}

\subsubsection{Лог}

\begin{lstlisting}[caption={Лог цели \texttt{report}},language={},captionpos=b]
Buildfile: /project/build.xml

report:

compile:
     [echo] Compiling the project from src/main/java...
     [echo] Source Java: 17
     [echo] Target Java: 17
   [delete] Deleting directory /project/build
    [mkdir] Created dir: /project/build
     [copy] Copying 2 files to /project/build
    [javac] Compiling 24 source files to /project/build
    [javac] warning: [options] system modules path not set in conjunction with -source 17
    [javac] 1 warning
     [echo] Done task 'compile'!

build:
     [echo] Building the project...
   [delete] Deleting directory /project/target
    [mkdir] Created dir: /project/target
     [copy] Copying 8 files to /project/target
     [copy] Copying 26 files to /project/target/WEB-INF/classes
     [copy] Copying 25 files to /project/target/WEB-INF/lib
      [jar] Building jar: /project/target/web3-1.0.0.war
     [echo] Done task 'build'!

test:
     [echo] Running JUnit tests...

compile-tests:
     [echo] Compiling JUnit tests...
   [delete] Deleting directory /project/build-test
    [mkdir] Created dir: /project/build-test
    [javac] Compiling 8 source files to /project/build-test
    [javac] warning: [options] system modules path not set in conjunction with -source 17
    [javac] 1 warning
     [echo] Done task 'compile-tests'!
     [java]
     [java] Thanks for using JUnit! Support its development at https://junit.org/sponsoring
     [java]
     [java] Test plan execution started. Number of static tests: 120
     [java] ╷
     [java] ├─ JUnit Jupiter
     [java] │  ├─ Тесты прямоугольника
     [java] │  │  ├─ Точка на нижней границе (y = 0) - должна вернуть true
     [java] │  │  │       tags: []
     [java] │  │  │   uniqueId: [engine:junit-jupiter]/[class:weblab.shapes.templates.RectangleTest]/[method:pointOnBottomBoundaryShouldReturnTrue()]
     [java] │  │  │     parent: [engine:junit-jupiter]/[class:weblab.shapes.templates.RectangleTest]
     [java] │  │  │     source: MethodSource [className = 'weblab.shapes.templates.RectangleTest', methodName = 'pointOnBottomBoundaryShouldReturnTrue', methodParameterTypes = '']
     [java] │  │  │   duration: 33 ms
     [java] │  │  │     status: ✔ SUCCESSFUL
     [java] │  │  ├─ Точка на правой границе (x = width) - должна вернуть true
     [java] │  │  │       tags: []
     [java] │  │  │   uniqueId: [engine:junit-jupiter]/[class:weblab.shapes.templates.RectangleTest]/[method:pointOnRightBoundaryShouldReturnTrue()]
     [java] │  │  │     parent: [engine:junit-jupiter]/[class:weblab.shapes.templates.RectangleTest]
     [java] │  │  │     source: MethodSource [className = 'weblab.shapes.templates.RectangleTest', methodName = 'pointOnRightBoundaryShouldReturnTrue', methodParameterTypes = '']
     [java] │  │  │   duration: 2 ms
     [java] │  │  │     status: ✔ SUCCESSFUL
     [java] │  │  ├─ Точка на левой границе (x = 0) - должна вернуть true
     [java] │  │  │       tags: []
     [java] │  │  │   uniqueId: [engine:junit-jupiter]/[class:weblab.shapes.templates.RectangleTest]/[method:pointOnLeftBoundaryShouldReturnTrue()]
     [java] │  │  │     parent: [engine:junit-jupiter]/[class:weblab.shapes.templates.RectangleTest]
     [java] │  │  │     source: MethodSource [className = 'weblab.shapes.templates.RectangleTest', methodName = 'pointOnLeftBoundaryShouldReturnTrue', methodParameterTypes = '']
     [java] │  │  │   duration: 2 ms
     [java] │  │  │     status: ✔ SUCCESSFUL

     <... cropped something>

     [java] │  │  ├─ Очистка пустой таблицы
     [java] │  │  │       tags: []
     [java] │  │  │   uniqueId: [engine:junit-jupiter]/[class:weblab.dao.ResultDAOTest]/[method:clearResultsEmptyTableShouldReturnZero()]
     [java] │  │  │     parent: [engine:junit-jupiter]/[class:weblab.dao.ResultDAOTest]
     [java] │  │  │     source: MethodSource [className = 'weblab.dao.ResultDAOTest', methodName = 'clearResultsEmptyTableShouldReturnZero', methodParameterTypes = '']
     [java] │  │  │   duration: 13 ms
     [java] │  │  │     status: ✔ SUCCESSFUL
     [java] │  └─ Тесты работы с БД finished after 3552 ms.
     [java] └─ JUnit Jupiter finished after 4899 ms.
     [java] Test plan execution finished. Number of all tests: 208
     [java]
     [java] Test run finished after 4974 ms
     [java] [        14 containers found      ]
     [java] [         0 containers skipped    ]
     [java] [        14 containers started    ]
     [java] [         0 containers aborted    ]
     [java] [        14 containers successful ]
     [java] [         0 containers failed     ]
     [java] [       208 tests found           ]
     [java] [         0 tests skipped         ]
     [java] [       208 tests started         ]
     [java] [         0 tests aborted         ]
     [java] [       208 tests successful      ]
     [java] [         0 tests failed          ]
     [java]
     [echo] Done task 'test'!
     [echo] Saving test reports to SVN...
     [java] Updating '.':
     [java] At revision 3.
     [java] Sending        build-test/reports/TEST-junit-jupiter.xml
     [java] Transmitting file data .done
     [java] Committing transaction...
     [java] Committed revision 4.
     [echo] Done task 'report'!

BUILD SUCCESSFUL
Total time: 28 seconds
\end{lstlisting}